\documentclass[11pt]{article}
\usepackage[utf8]{inputenc}
\usepackage{amsmath,amssymb}
\usepackage{graphicx}
\usepackage{booktabs}
\usepackage{authblk}
\usepackage[margin=1in]{geometry}
\usepackage[hidelinks]{hyperref}
\usepackage{xurl}
\usepackage{setspace}

\title{The Binding Energy Problem QCD Never Solved\\ --- and a Coulomb Model That Does}
\author[1]{Claes Johnson}
\affil[1]{KTH Royal Institute of Technology, SE-100 44 Stockholm, Sweden \\
\texttt{claesjohnson@gmail.com}}
\date{\today \\[0.4em]
\small\textit{Argument and claims by the author; drafting and the cited computations prepared with AI assistance (Claude, Anthropic), reviewed and verified by the author.}}

\begin{document}
\maketitle

\begin{abstract}
Quantum Chromodynamics (QCD) was put forward in 1973 as the fundamental theory of the strong
interaction, and with it the accepted account of why atomic nuclei exist and hold together. Fifty-three
years later, QCD has not produced a single \emph{parameter-free, first-principles} prediction of a
nuclear binding energy --- not the deuteron, not the alpha particle. This is not a claim that QCD is
wrong: its domain of demonstrated success is the perturbative, short-distance regime of quarks and
gluons, where asymptotic freedom makes it tractable. But nuclear binding lives in the opposite,
long-distance, non-perturbative regime, and there the binding energies have in practice been supplied not
by QCD but by \emph{fitted} effective theories, whose low-energy constants are read off the very nuclear
data they then reproduce; lattice QCD, meanwhile, reaches light nuclei only at unphysically large quark
masses, with uncontrolled extrapolation to the physical point. We document this gap, distinguish it
carefully from what QCD \emph{has} established, and then set beside it an alternative: the RealNucleus
model, in which a nucleus is protons and electrons as Coulomb charge clouds, bound by the electric force
alone --- no strong force, no weak force. With a single scale calibrated on the deuteron and nothing else
fitted, this model reproduces the alpha binding energy ($\sim 28$~MeV, $103\%$ of the measured value), the
alpha-to-deuteron binding ratio ($13.1$ against a measured $12.7$, a scale-free prediction), the
alpha-conjugate binding ladder up to ${}^{40}$Ca at $\sim 107$--$108\%$, and the emergence of saturation. The contrast is stark: the theory
written to explain nuclear binding has never delivered it parameter-free, while a model containing no
strong force at all delivers it with one scale. We argue this makes a long-avoided question legitimate
rather than settled.

\medskip
\noindent\textbf{Keywords:} nuclear binding energy; QCD; strong force; lattice QCD; chiral effective
field theory; deuteron; alpha particle; RealNucleus; Coulomb.
\end{abstract}

\doublespacing

\section{The question}
\label{sec:question}

Why does a nucleus hold together? In the standard picture a deuteron is a single proton bound to a
neutron, and an alpha two protons and two neutrons. The deuteron thus carries only one proton, so no
Coulomb repulsion acts within it: its $2.224$~MeV binding is purely the strong proton--neutron
attraction. Already in the alpha, however, the two protons repel one another electrically and the nucleus
holds nonetheless. Something must therefore both bind the nucleons together and, wherever protons sit
side by side, overcome their Coulomb repulsion, supplying the observed binding energy: $2.224$~MeV for
the deuteron, $28.3$~MeV for the alpha, rising to $\sim 8.8$~MeV per nucleon near iron.

It is worth isolating the deuteron, for it poses the question in its purest form. With only one proton
there is no Coulomb repulsion to overcome; the $2.224$~MeV is the strong nucleon--nucleon attraction
acting alone, unobscured by electrostatics or by the many-body complications of larger nuclei. What
binds a single proton to a single neutron? This is the simplest bound state in nuclear physics and the
first number a fundamental theory of the strong force ought to deliver. As we record below, it has not
been delivered parameter-free either --- so the shortfall cannot be blamed on the difficulty of the
multi-proton Coulomb problem; it is present already where no Coulomb problem exists.

The standard answer is the strong nuclear force, and its fundamental theory is Quantum Chromodynamics.
The purpose of this paper is narrow and, we believe, factual: to record that after fifty-three years this
theory has not predicted, from first principles and without fitting, the binding energy of any nucleus;
to state precisely what that does and does not mean; and to place beside it a model that reproduces those
binding energies from the electric force alone. We are not attacking QCD on its own ground. We are asking
whether the ground it was supposed to cover has in fact been covered.

\section{What QCD was meant to do for nuclei}
\label{sec:mandate}

QCD was formulated in 1973 with the discovery of asymptotic freedom
\cite{gross1973,politzer1973}: the coupling between quarks and gluons weakens at short distance and
strengthens at long distance. This single structural fact organizes everything that follows. Where the
coupling is weak --- high momentum transfer, short distance --- QCD is calculable and has been tested
extensively. Where the coupling is strong --- low energy, long distance, the scale of confinement and of
whole nucleons bound into nuclei --- it is not perturbatively calculable at all.

This is the essential point, and it is not a criticism levelled from outside: \emph{a nucleus is a
low-energy, long-distance object.} Nucleons a femtometre apart, bound by an MeV-scale energy, sit
squarely in the regime where QCD's own coupling is large and its own perturbation theory fails. The
theory is at its most powerful precisely where nuclei are not, and at its weakest precisely where they
are. Whatever QCD has demonstrated in the perturbative ultraviolet, the existence and binding of nuclei
is a question in the non-perturbative infrared, and it is there that the account must be delivered.

Two routes have been pursued to deliver it. Neither has produced a parameter-free binding energy, and it
is worth being exact about why.

\section{Route one: lattice QCD}
\label{sec:lattice}

Lattice QCD evaluates the theory numerically on a discretized spacetime, and is the only route that is
QCD \emph{as such}, with no additional physical assumptions. For single hadrons --- the proton mass, the
low-lying meson and baryon spectrum --- it has succeeded at the percent level, and this is a genuine
achievement of the theory in its strong-coupling regime.

For \emph{nuclei} the situation is entirely different. Multi-nucleon systems on the lattice are afflicted
by a signal-to-noise problem that grows exponentially with the number of quarks, and by contamination
from the many nearby scattering states. To make the computation tractable at all, the existing
calculations of light nuclei --- the deuteron, ${}^3$He, ${}^4$He --- have been carried out at
\emph{unphysically large} quark masses, corresponding to pion masses several times the physical value
(typically in the range $m_\pi \sim 300$--$800$~MeV against the physical $m_\pi \simeq 140$~MeV)
\cite{beane2013,yamazaki2012}. At those heavy masses a bound deuteron and alpha can be seen, but the
numbers are not the physical binding energies; the extrapolation down to the physical point is long,
uncontrolled, and openly disputed, and different collaborations have not converged even on whether the
two-nucleon system is bound at the physical point.

The honest summary is this: fifty-three years after QCD, there exists no lattice calculation that
predicts the measured binding energy of the deuteron, or of the alpha, at the physical quark masses to
controlled accuracy. The one route that is QCD without further assumptions has not reached a single real
nucleus.

\section{Route two: effective field theory --- and the quiet handoff}
\label{sec:eft}

The nuclear binding energies that theory does reproduce are obtained by a different route: chiral
effective field theory (\(\chi\)EFT) and phenomenological nucleon--nucleon potentials
\cite{weinberg1990,epelbaum2009}. Here one does not solve QCD. One writes the most general interaction
among nucleons and pions consistent with the symmetries of QCD, organized as an expansion in
low momentum, and then \emph{fixes the coefficients of that expansion --- the low-energy constants --- by
fitting to nuclear and nucleon--nucleon scattering data.} With enough such constants determined from
experiment, the light nuclei and their binding energies are reproduced, often well.

This is useful and it is legitimate physics. But it must be stated plainly what it is and is not. The
binding energies so obtained are not \emph{predictions from QCD}; the low-energy constants are read off
the same nuclear world whose binding they then describe. The values of those constants are not, at
present, computed from QCD. In effect the problem of nuclear binding has been quietly handed off from the
fundamental theory to a fitted parametrization --- a parametrization inspired by QCD's symmetries, but
calibrated on the answer. The strong force explains the existence of nuclei in the same sense that a
polynomial fitted through data points ``explains'' the data: it accommodates them; it does not predict
them from below.

\section{The gap, stated precisely}
\label{sec:gap}

It is important to claim neither too much nor too little.

\emph{We do not claim} that QCD is falsified, that it fails on quarks and gluons, or that its
strong-coupling regime is understood to be empty. Its successes on the hadron spectrum are real, and
belong to the regime where it is calculable.

\emph{We do claim}, and take to be simply factual, that after fifty-three years:
\begin{itemize}
\item no calculation of QCD \emph{as such} (lattice, at physical quark masses, with controlled
  extrapolation) predicts the binding energy of any nucleus, the deuteron and alpha included;
\item the binding energies that theory reproduces come from effective theories whose parameters are
  fitted to nuclear data, not derived from QCD;
\item hence the specific number that motivated the strong force in the first place --- the energy that
  binds a nucleus against Coulomb repulsion --- is still not among the numbers the fundamental theory
  predicts.
\end{itemize}

That a theory can be enormously successful in one regime and, for half a century, silent on the very
phenomenon it was built to explain in another, is worth stating aloud rather than passing over.

\section{A model that delivers the binding energies: RealNucleus}
\label{sec:realnucleus}

Set beside this gap a model of an entirely different character. In the RealNucleus
picture \cite{realnucleus} there is no strong force and no weak force. A nucleus is nothing but
\emph{protons and electrons as charge clouds}, bound by the ordinary Coulomb attraction --- the same
electric law that binds atoms and molecules, read with the charges rearranged. The neutron is a bound
proton--electron pair; the deuteron is $2p+1e$, two positive charges glued by one negative one, the
nuclear cousin of the molecular ion $\mathrm{H}_2^+$. The energy minimized is the plain Coulomb
functional

\paragraph{Where the two pictures already part: the deuteron.} The contrast is sharpest at the very
first nucleus. In the standard picture the deuteron is a proton bound to a \emph{neutron}, and the open
question --- unanswered parameter-free by QCD, as recorded above --- is what binds them: what supplies
the $2.224$~MeV attraction between a proton and a neutral neutron. RealNucleus does not answer that
question; it dissolves it. There is no elementary neutron to bind: the neutron is itself a bound
$p+e$ pair, so the deuteron is not proton-plus-neutron but $2p+1e$ --- two like charges that
\emph{repel}, held together by one shared electron, exactly as one electron binds two protons in
$\mathrm{H}_2^+$. The binding is then no longer a mysterious strong attraction between a charge and a
neutral object, but ordinary electrostatics: the electron's attraction to both protons outweighs the
proton--proton repulsion. The two theories thus disagree about the deuteron already at the level of what
it is made of --- proton$+$neutron versus $2p+1e$ --- and about what holds it together --- an
undelivered strong force versus the Coulomb law that also binds $\mathrm{H}_2^+$.

Concretely, the energy minimized is the plain Coulomb functional
\[
  E \;=\; \sum_i \tfrac12\!\int |\nabla \psi_i|^2
        \;+\; \tfrac12\!\iint \frac{\rho(x)\,\rho(y)}{|x-y|}\,,
  \qquad \rho = \sum_i q_i \psi_i^2,\quad q_i = \pm 1,
\]
over non-overlapping charge domains with free boundaries.

From this, with the \emph{electric force only} and a \emph{single scale} fixed on the deuteron ---
nothing else fitted --- the model delivers:
\begin{itemize}
\item the \textbf{alpha binding energy}, $\sim 28$~MeV (computed $29.1$~MeV against the measured $28.3$,
  $103\%$) --- the very number the two QCD routes do not give parameter-free;
\item the \textbf{alpha-to-deuteron binding ratio}, $13.1$ against a measured $12.7$ --- a genuinely
  \emph{scale-free} prediction, since a ratio does not see the calibration;
\item the \textbf{alpha-conjugate binding ladder}, ${}^4$He, ${}^{12}$C, ${}^{16}$O, \dots, ${}^{40}$Ca,
  reproduced at $\sim 107$--$108\%$ of the measured binding, with near-constant binding per nucleon
  (\textbf{saturation}) emerging rather than imposed;
\item the fusion reaction $d+d\to{}^4$He and alpha decay (Gamow / Geiger--Nuttall), from the same
  functional;
\item and a result that the \textbf{electron's mass is irrelevant} to the structure: a genuinely light
  electron, allowed to relax, chooses a flat, charge-continuous shape, so the nuclear scale is set by the
  heavy proton and the atomic scale by the light electron --- two sizes, one Coulomb law.
\end{itemize}

\paragraph{The honest caveat.} This is one scale, not literally zero input: the deuteron energy sets the
unit. But a unit is not a fit. Once it is chosen, every \emph{ratio} and the \emph{shape} of the
binding-per-nucleon curve are predictions, not adjustments --- in contrast to the many fitted low-energy
constants of the effective-theory route. There are genuine open problems: the spin--statistics of the
electron confined in the nucleus, the closed-shell structure, and the model stays deliberately silent on
the neutrino. None of this is settled, and it is offered so that it can be checked and refuted.

\section{The contrast}
\label{sec:contrast}

The two accounts can be placed side by side without rhetoric.

\begin{center}
\begin{tabular}{@{}p{0.46\linewidth} p{0.46\linewidth}@{}}
\toprule
\textbf{QCD (strong force), since 1973} & \textbf{RealNucleus (Coulomb alone)} \\
\midrule
Fundamental theory of the strong interaction; the accepted reason nuclei exist. &
No strong force, no weak force; nuclei bound by the electric force only. \\[4pt]
Calculable in the perturbative UV; nuclear binding lives in the non-perturbative IR. &
The same Coulomb functional at all scales; two sizes (proton, electron) from one law. \\[4pt]
Lattice QCD (the theory as such) reaches light nuclei only at unphysical quark masses; no controlled
binding energy at the physical point. &
Alpha $\sim 28$~MeV ($103\%$) and the ladder to ${}^{40}$Ca at $\sim 107$--$108\%$ from one deuteron-fixed scale. \\[4pt]
$\chi$EFT reproduces nuclei only with low-energy constants \emph{fitted} to nuclear data. &
No fitted constants beyond the single overall scale; ratios and saturation are scale-free. \\
\bottomrule
\end{tabular}
\end{center}

The alpha particle's $\sim 28$~MeV is the canonical thing the strong force was invented to account for.
It is reproduced, to about $103\%$, with a single scale, by a model that contains no strong force at all.

\section{Discussion}
\label{sec:discussion}

None of this retires QCD as the theory of quarks and gluons. What it does is make a question legitimate
that the field has grown used to leaving aside: if the binding energy of a nucleus can be obtained from
the electric force with one calibration constant, how much of the strong-force machinery is actually
required to explain why nuclei exist, and how much of what passes for explanation is fitting?

The reasonable objection --- ``QCD is correct, the nuclear regime is merely hard'' --- may well be true.
But ``merely hard'' has meant fifty-three years without a single parameter-free nuclear binding energy,
while the binding was supplied by fitted parametrizations and, separately, reproduced from scratch by a
model with no strong force in it. A theory earns its standing by predicting the phenomena it was built
for. On the specific phenomenon of nuclear binding, that debt is still outstanding, and the existence of
a Coulomb-only account that does settle it --- with its own open problems, offered for refutation ---
is reason to ask the question out loud rather than to assume it has already been answered.

\begin{thebibliography}{99}

\bibitem{gross1973} D.~J.~Gross and F.~Wilczek,
\emph{Ultraviolet Behavior of Non-Abelian Gauge Theories},
Phys. Rev. Lett. \textbf{30}, 1343 (1973).

\bibitem{politzer1973} H.~D.~Politzer,
\emph{Reliable Perturbative Results for Strong Interactions?},
Phys. Rev. Lett. \textbf{30}, 1346 (1973).

\bibitem{beane2013} S.~R.~Beane \emph{et al.} (NPLQCD Collaboration),
\emph{Light nuclei and hypernuclei from quantum chromodynamics in the limit of SU(3) flavor symmetry},
Phys. Rev. D \textbf{87}, 034506 (2013).

\bibitem{yamazaki2012} T.~Yamazaki \emph{et al.},
\emph{Helium nuclei in quenched and (2+1)-flavor lattice QCD}, and related PACS-CS studies,
Phys. Rev. D \textbf{86}, 074514 (2012).

\bibitem{weinberg1990} S.~Weinberg,
\emph{Nuclear forces from chiral Lagrangians},
Phys. Lett. B \textbf{251}, 288 (1990).

\bibitem{epelbaum2009} E.~Epelbaum, H.-W.~Hammer, and Ulf-G.~Mei\ss ner,
\emph{Modern theory of nuclear forces},
Rev. Mod. Phys. \textbf{81}, 1773 (2009).

\bibitem{realnucleus} C.~Johnson,
\emph{RealNucleus: nuclei by Coulomb forces alone},
manuscript (2026); simulations at \url{https://claes542.github.io/RealMolecule/gallery.html}.

\end{thebibliography}

\end{document}
